% NOETHER PAPER 2 — KOREAN TRANSLATION-PRODUCER DRAFT — UNCHECKED
% Unit: T13-U70; current-authority whole lines 1726--1744
% Source slice: 2,091 LF UTF-8 bytes; SHA-256 A4EE5C15092F16E7BD8BDCB8A851BEFBC7313CBCACB0A082AA0B41BF62A2F0E5
% Pointer: NOETH-DE-AUTH-v005-20260804; SHA-256 42E6844BFCBFB2133E9AA323A823604351CF9C49550AFCF34ECAAF7887185660
% This is producer translation only: no source/Korean/formula review, compilation, or rendering.

c) 형식열 $\theta_\eta^3$: 형식열 $\theta_\eta H_\vartheta$와 $\theta_\eta^2H_\vartheta$에 동시에 속하는 형식들, 즉 형식열 $\theta_\eta^2H_\vartheta$의 형식들을 이중 환원하여 얻는다. 한편으로는 더 높은 기호를 갖는 형식들로 환원하고, 다른 한편으로는 더 높은 기호를 갖는 형식들과 더 높은 형식열 $\theta_\eta^3$으로 환원한다.

d) 형식 $\theta_\eta^2(\vartheta\eta x)$, $\theta_\eta^2(\vartheta\eta x)^2$, $\theta_\eta^3(\vartheta\eta x)$: \S\ 9의 계산과 마찬가지로 환원자 $a$를 이용한 이중 환원으로 얻는다.

e) 형식 $\theta_\eta^2(\theta Hu)^2$와 $\theta_\eta^2(\vartheta\eta x)^2$: 환원자 $\A_\nu^2$와 $(a\A u)^4$를 이용하여 식 $\A_\nu^2(a\A u)^4$와 $(a\A\widehat{\nu}x)^4$을 이중 환원함으로써 얻는다.

f) 형식 $H_\vartheta$와 $H_\vartheta^2$: 분해 가능 형식이다. $H_\vartheta^2$의 경우 이는 식 $\A_\nu^2(a\A u)^2$의 이중 환원으로, 또는 곱 $(a_\nu^2)^2$, $a_\nu^2\cdot b_{\nu_1}$을 계의 형식들로 직접 계산함으로써 얻는다. ($(a_\nu)^2$을 이와 같이 계산하면 분해 가능 형식들 사이의 관계가 생긴다.)

g) 형식열 $\theta\cdot H$의 형식들이 두 가지 환원 방법을 허용하기 때문에, 즉 두 환원 가능한 형식열에 속하기 때문에 다음과 같이 더 높은 형식들이 환원된다.
\[
\begin{array}{rl}
g&\text{는 }\theta_\eta^2(\vartheta\eta x)^2\text{의 이중 환원으로: 환원 d)와 e)가 가능하다;}\\[0.3em]
s_\vartheta^2(\theta s u)&\text{는 }\theta_\eta^3(\vartheta\eta x)\text{의 이중 환원으로: 환원 c)와 d)가 가능하다;}\\[0.3em]
s_\vartheta^2(\theta s u)^2&\text{은 }\theta_\eta^2H_\vartheta^2\text{의 이중 환원으로: 환원 a)가 가능하며}\\
&\text{환원자 }\theta_\nu^2(\vartheta\nu x)^2\text{을 인수로 갖는다;}\\[0.3em]
s_\nu^2&\text{은 }\theta_\eta^3H_\vartheta\text{의 이중 환원으로: 환원 a)와 c)가 가능하다.}
\end{array}
\]
나머지 형식들의 독립성은 형식열 $\A_\nu^2(a\A u)^2\equiv\theta_\eta^2$의 이중 환원으로 증명된다.
